\documentclass[11pt,a4paper]{article}

\usepackage[utf8]{inputenc}
\usepackage[T1]{fontenc}
\usepackage[brazilian]{babel}
\usepackage{graphicx}
\usepackage{booktabs}
\usepackage{array}
\usepackage{float}
\usepackage{geometry}
\usepackage{longtable}
\usepackage{url}
\usepackage{amsmath}
% Permite quebrar texto monoespaçado (domínios) após pontos e barras.
\renewcommand{\UrlBreaks}{\do\.\do\/\do\-\do\_}
\usepackage[hidelinks]{hyperref}

\geometry{top=2.5cm, bottom=2.5cm, left=2.5cm, right=2.5cm}

% Figura com largura relativa e altura máxima, mantendo a proporção.
\newcommand{\evid}[3]{%
  \begin{figure}[H]
    \centering
    \includegraphics[width=#2\textwidth,height=0.72\textheight,keepaspectratio]{#1}
    \caption{#3}
  \end{figure}}

% Três capturas do popup lado a lado.
\newcommand{\popups}[4]{%
  \begin{figure}[H]
    \centering
    \includegraphics[width=0.31\textwidth]{#1}\hfill
    \includegraphics[width=0.31\textwidth]{#2}\hfill
    \includegraphics[width=0.31\textwidth]{#3}
    \caption{#4}
  \end{figure}}

% Duas capturas lado a lado, com a mesma altura.
\newcommand{\pair}[3]{%
  \begin{figure}[H]
    \centering
    \includegraphics[height=6.2cm,width=0.66\textwidth,keepaspectratio]{#1}\hfill
    \includegraphics[height=6.2cm,width=0.3\textwidth,keepaspectratio]{#2}
    \caption{#3}
  \end{figure}}

% Duas capturas empilhadas, cada uma na largura total.
\newcommand{\stack}[3]{%
  \begin{figure}[H]
    \centering
    \includegraphics[width=\textwidth,height=0.42\textheight,keepaspectratio]{#1}\\[4pt]
    \includegraphics[width=\textwidth,height=0.42\textheight,keepaspectratio]{#2}
    \caption{#3}
  \end{figure}}

\newcommand{\ok}{sim}
\newcommand{\no}{não}

\title{
    Privacy Inspector:\\
    Detecção e Análise de Rastreadores e Mecanismos de Armazenamento no Firefox
}
\author{Cynthia Naoko Yasutake\\Insper}
\date{Setembro de 2026}

\begin{document}
\sloppy

\maketitle

\section{Introdução}

Este relatório apresenta o desenvolvimento e a validação do Privacy Inspector, uma extensão para o Mozilla Firefox que identifica, por aba, conexões com domínios de terceiros, cookies definidos durante o carregamento, uso de armazenamento HTML5, canvas fingerprinting, bounce tracking e cookie sync, indicadores de sequestro de navegador (hijacking e hook) e atribui à página uma pontuação de privacidade com metodologia explícita. A extensão também permite ao usuário manter uma lista de bloqueio personalizada.

A validação foi feita em duas frentes: nas DuckDuckGo Privacy Test Pages, que oferecem cenários padronizados com resultado esperado conhecido, e em três sites reais (G1, University of British Columbia e YouTube), cujos resultados foram comparados com os arquivos HAR exportados do DevTools, com o relatório do Blacklight (The Markup) e com os bloqueios do uBlock Origin. O código, as instruções de instalação e todas as evidências estão no repositório, na pasta \path{privacy-inspector/relatorio/evidencias/}.

\section{Arquitetura e metodologia do plugin}

O Privacy Inspector é uma WebExtension em Manifest V3 composta pelos seguintes arquivos:

\begin{itemize}
  \item \texttt{background.\allowbreak{}js}: monitora a rede com a API \texttt{webRequest} e a navegação com a API \texttt{webNavigation}, mantém os dados de cada aba em \texttt{storage.\allowbreak{}session} e aplica a lista de bloqueio;
  \item \texttt{content-script.\allowbreak{}js}: executa no contexto isolado da extensão, inspeciona o armazenamento HTML5 e repassa ao background as mensagens dos detectores que rodam no contexto da página;
  \item \texttt{canvas-detector.\allowbreak{}js} e \texttt{hook-detector.\allowbreak{}js}: executam no contexto da própria página (\texttt{world: MAIN}) em \texttt{document\_start}, antes dos scripts do site;
  \item \texttt{score.\allowbreak{}js}: implementa a pontuação de privacidade;
  \item \texttt{popup.\allowbreak{}html}, \texttt{popup.\allowbreak{}js} e \texttt{popup.\allowbreak{}css}: apresentam o relatório da aba ativa e a lista de bloqueio.
\end{itemize}

\subsection{Domínios de terceiros}

A cada novo carregamento do documento principal, o plugin registra o domínio registrável da página (eTLD+1, obtido com a API \texttt{publicSuffix}). Toda requisição subsequente da aba cujo domínio registrável difere do da página é contada como de terceira parte. Como a lista de sufixos públicos inclui domínios privados como \texttt{googleapis.\allowbreak{}com}, subdomínios como \texttt{fonts.\allowbreak{}googleapis.\allowbreak{}com} e \texttt{ajax.\allowbreak{}googleapis.\allowbreak{}com} aparecem como domínios registráveis distintos.

\subsection{Cookies}

Os cookies são identificados pelos cabeçalhos \texttt{Set-Cookie} recebidos em \path{webRequest.onHeadersReceived}. São de primeira parte quando a resposta pertence ao domínio da página e de terceira parte caso contrário; são persistentes quando possuem \texttt{Expires} ou \texttt{Max-Age} e de sessão caso contrário. A contagem representa eventos de definição, de modo que um mesmo cookie redefinido várias vezes é contado mais de uma vez. Cookies criados por \texttt{document.\allowbreak{}cookie} não são contabilizados.

\subsection{Armazenamento HTML5}

O content script conta os itens de \texttt{localStorage} e \texttt{sessionStorage} e os bancos retornados por \texttt{indexedDB.\allowbreak{}databases()} no documento principal, cerca de um segundo após o início do carregamento.

\subsection{Canvas fingerprinting}

O \texttt{canvas-detector.\allowbreak{}js} substitui, no protótipo, os métodos de leitura de canvas usados para gerar impressões digitais: \texttt{HTMLCanvasElement.\allowbreak{}toDataURL}, \texttt{HTMLCanvasElement.\allowbreak{}toBlob}, \texttt{CanvasRenderingContext2D.\allowbreak{}getImageData}, \texttt{OffscreenCanvas.\allowbreak{}convertToBlob} e \texttt{OffscreenCanvas.\allowbreak{}transferToImageBitmap}. Cada chamada é repassada à função original e registrada por método.

\subsection{Bounce tracking, cookie sync e parâmetros de rastreamento}

Três mecanismos são observados:
\begin{itemize}
  \item \textbf{Bounce tracking}: uma navegação A $\rightarrow$ X $\rightarrow$ B em que X é outro site em relação a A, permanece aberto por menos de 10 segundos e sai de forma automática, seja por redirecionamento HTTP, seja por JavaScript (indicado pelo \texttt{transitionQualifier} do Firefox ou por uma permanência inferior a 2 segundos). Quando B é outro subdomínio do mesmo site que X, exige-se que a URL de B carregue um parâmetro com aparência de identificador, para não confundir o bounce com redirecionamentos comuns como \texttt{example.\allowbreak{}com} $\rightarrow$ \texttt{www.\allowbreak{}example.\allowbreak{}com};
  \item \textbf{Cookie sync}: um redirecionamento de sub-recurso (em geral um pixel) entre dois domínios de terceiros diferentes, levando um identificador na URL;
  \item \textbf{Parâmetros de rastreamento} na URL da página, como \texttt{utm\_*}, \texttt{fbclid}, \texttt{gclid} e \texttt{msclkid}.
\end{itemize}

\subsection{Indicadores de hijacking e hook}

O plugin registra como indicadores de hijacking: WebSocket aberto para domínio de terceiro; polling persistente, definido como cinco ou mais requisições XHR/fetch para o mesmo terceiro após os primeiros 10 segundos da página; abertura de nova aba ou janela pela página para outro domínio (\texttt{window.\allowbreak{}open} e \path{webNavigation.onCreatedNavigationTarget}), que cobre casos de pop-under; redirecionamento automático para outro domínio sem o padrão de bounce; e sobrescrita de objetos globais. Para este último, o \texttt{hook-detector.\allowbreak{}js} guarda as referências nativas de \texttt{fetch}, \texttt{XMLHttpRequest.\allowbreak{}prototype.\allowbreak{}open/\allowbreak{}send}, \texttt{WebSocket}, \texttt{navigator.\allowbreak{}sendBeacon}, \texttt{document.\allowbreak{}write}, \texttt{EventTarget.\allowbreak{}prototype.\allowbreak{}addEventListener} e \texttt{window.\allowbreak{}open} antes de qualquer script da página e as compara 1, 5, 15 e 30 segundos após o evento \texttt{load}. Esses são os pontos de interceptação usados por frameworks de hook como o BeEF.

\subsection{Lista de bloqueio personalizada}

O usuário adiciona domínios no popup, e a lista é salva em \texttt{storage.\allowbreak{}local}. O listener de \path{webRequest.onBeforeRequest}, registrado como \texttt{blocking} com a permissão \texttt{webRequestBlocking}, cancela requisições para os domínios da lista e seus subdomínios, exceto o documento principal. As requisições canceladas aparecem no popup por domínio.

\subsection{Condições de coleta}

As coletas finais foram feitas em 28/09/2026, no Firefox 156 para Windows, em janela privada nova para cada site, sem aceitar banners de consentimento e com o uBlock Origin desativado. A proteção contra rastreamento do Firefox (ETP) permaneceu na configuração padrão da janela privada, que bloqueia conteúdo de rastreamento conhecido; seus bloqueios aparecem nos HARs como requisições com status 0 e são considerados na análise. O uBlock Origin foi ativado apenas nas coletas específicas do seu logger. O Blacklight foi executado na versão desktop, a partir de servidores na Califórnia (parâmetro \texttt{location=us-ca}).

Os HARs publicados foram sanitizados: valores de cookies, cabeçalhos \texttt{Cookie} e \texttt{Referer}, corpos de requisição e de resposta e parâmetros de URL com aparência de identificador foram substituídos por \texttt{REDACTED}, preservando domínios, caminhos, métodos, códigos de status e nomes de cookies. Os HARs brutos não são versionados. Os arquivos \texttt{*-etp.\allowbreak{}sanitized.\allowbreak{}har} correspondem à coleta inicial, de 27/09, feita em janela normal e com uma versão anterior do plugin.

\section{DuckDuckGo Privacy Test Pages}

Cada página foi aberta e, nos testes interativos, os passos indicados foram executados antes de abrir o popup. Nos testes com atraso, aguardou-se o tempo indicado pela página. As tabelas comparam o resultado reportado pela página com o observado no plugin e explicam cada divergência.

\subsection{Tracker Reporting}

\small
\begin{longtable}{|p{2.6cm}|p{3cm}|p{3.4cm}|p{5.0cm}|}
\caption{Resultados de Tracker Reporting.}\\
\hline
\textbf{Teste} & \textbf{Esperado (página)} & \textbf{Plugin} & \textbf{Explicação} \\ \hline
\endfirsthead
\hline
\textbf{Teste} & \textbf{Esperado (página)} & \textbf{Plugin} & \textbf{Explicação} \\ \hline
\endhead
Major tracker via script & 1 major tracker carregado por \texttt{script src} & \texttt{doubleclick.\allowbreak{}net} como terceiro; 0 cookies; score 99 & A requisição foi observada em \texttt{onBeforeRequest}. O script do \texttt{doubleclick.\allowbreak{}net} não definiu cookies na resposta, e o plugin identifica o domínio como terceiro, mas não o classifica como rastreador por lista. \\ \hline
Major tracker com surrogate & 1 major tracker substituível por surrogate & \texttt{doubleclick.\allowbreak{}net}; 0 cookies; score 99 & O plugin detecta a conexão, mas não substitui scripts por surrogates. No logger do uBlock Origin do G1, os scripts \texttt{adsbygoogle.\allowbreak{}js} e \texttt{gpt.\allowbreak{}js} foram trocados por versões neutralizadas (\texttt{redirect-rule}), mecanismo que este teste avalia e que o plugin não implementa. \\ \hline
Major tracker via img & 1 major tracker carregado por \texttt{img src} & \texttt{facebook.\allowbreak{}com}; 1 cookie de terceira parte persistente; score 96 & Detectou a conexão e o \texttt{Set-Cookie} da resposta da imagem, penalizado no critério de cookies de terceira parte. \\ \hline
Tracker via DocumentFragment & Requisição criada dinamicamente por fragmento & \texttt{facebook.\allowbreak{}com}; 1 cookie de terceira parte de sessão; score 97 & Detectado porque o monitoramento é feito na rede, independente de o elemento existir no HTML original. O plugin não informa que a origem foi um \texttt{DocumentFragment}. \\ \hline
DocumentFragment com atraso de 5 s & Mesma requisição, 5 s após o carregamento & \texttt{facebook.\allowbreak{}com}; 1 cookie de terceira parte de sessão; score 97 & O monitoramento continua ativo depois do carregamento inicial, então a requisição tardia também foi registrada. \\ \hline
Major tracker via fetch & Tracker carregado com a API Fetch & \texttt{facebook.\allowbreak{}com}; 1 cookie de terceira parte persistente; score 96 & Detectado na rede. O tipo da requisição (fetch) não é exibido no popup. \\ \hline
Fetch com atraso de 5 s & Mesma requisição, 5 s após o carregamento & \texttt{facebook.\allowbreak{}com}; 1 cookie de terceira parte persistente; score 96 & Detectado mesmo após o atraso. \\ \hline
\end{longtable}
\normalsize

Em todos os casos o plugin registrou 1 item em \texttt{localStorage}, 1 em \texttt{sessionStorage} e 1 banco IndexedDB, criados pela própria página de testes. O ícone do uBlock Origin não aparece nas capturas porque a extensão estava desativada.

\stack{evidencias/ddg/tracker-reporting/1-major-via-script-crop.png}{evidencias/ddg/tracker-reporting/1-major-via-surrogate-crop.png}{Major tracker via script (acima) e com surrogate (abaixo).}
\stack{evidencias/ddg/tracker-reporting/1-major-via-img-crop.png}{evidencias/ddg/tracker-reporting/1-major-via-img-1.png}{Major tracker via img: popup e seção de cookies.}
\stack{evidencias/ddg/tracker-reporting/tracker-document-fragment-crop.png}{evidencias/ddg/tracker-reporting/tracker-document-fragment-1.png}{Tracker via DocumentFragment.}
\stack{evidencias/ddg/tracker-reporting/tracker-document-fragment-delay-5s-crop.png}{evidencias/ddg/tracker-reporting/tracker-document-fragment-delay-5s-1-crop.png}{Tracker via DocumentFragment com atraso de 5 s.}
\stack{evidencias/ddg/tracker-reporting/1-major-via-fetch-crop.png}{evidencias/ddg/tracker-reporting/1-major-via-fetch-1-crop.png}{Major tracker via fetch.}
\stack{evidencias/ddg/tracker-reporting/1-major-via-fetch-delay-5s-crop.png}{evidencias/ddg/tracker-reporting/1-major-via-fetch-delay-5s-1-crop.png}{Major tracker via fetch com atraso de 5 s.}

\subsection{Storage Blocking}

\begin{table}[H]
\centering\small
\begin{tabular}{|p{2.4cm}|p{3.2cm}|p{3.6cm}|p{4.8cm}|}
\hline
\textbf{Teste} & \textbf{Esperado (página)} & \textbf{Plugin} & \textbf{Explicação} \\ \hline
Storage Blocking & Armazenar e recuperar um valor por 23 mecanismos; a página reportou duas falhas & 1 item em \texttt{localStorage}, 1 em \texttt{sessionStorage}, 1 banco IndexedDB, 2 domínios de terceiros e 9 cookies & O plugin confirmou os três mecanismos HTML5 do documento principal, mas não inspeciona o armazenamento dos iframes de terceiros nem mecanismos como Cache API, CookieStore, Service Worker e \texttt{window.\allowbreak{}name}. A contagem de cookies considera apenas \texttt{Set-Cookie}, enquanto a página também cria cookies por JavaScript. \\ \hline
\end{tabular}
\caption{Resultado de Storage Blocking.}
\end{table}

\evid{evidencias/ddg/storage-blocking/storage-blocking-crop.png}{1}{Storage Blocking (coleta de 27/09, versão anterior do popup).}

\subsection{Fingerprinting e Canvas}

\begin{table}[H]
\centering\small
\begin{tabular}{|p{2.4cm}|p{3.2cm}|p{3.6cm}|p{4.8cm}|}
\hline
\textbf{Teste} & \textbf{Esperado (página)} & \textbf{Plugin} & \textbf{Explicação} \\ \hline
Fingerprinting geral & Coleta de 124 características do navegador (14 falharam) & Canvas detectado: sim; 6 chamadas & O plugin identificou a parte da coleta feita por canvas: \texttt{toDataURL} (4), \texttt{getImageData} (1) e \texttt{OffscreenCanvas.\allowbreak{}transferToImageBitmap} (1). Os dados obtidos de \texttt{navigator}, \texttt{screen}, WebGL e fontes não são monitorados e, portanto, não aparecem no popup. \\ \hline
Canvas visual check & Renderização e leitura por data URL, \texttt{getImageData} e \texttt{OffscreenCanvas} & Mesmas 6 chamadas, por método & As quatro formas de leitura exibidas pela página correspondem aos métodos interceptados. A imagem é gerada normalmente, pois o plugin apenas observa e não altera o retorno. \\ \hline
\end{tabular}
\caption{Resultados de fingerprinting.}
\end{table}

\evid{evidencias/ddg/fingerprinting-canvas/fingerprinting-general-crop.png}{1}{Fingerprinting geral: canvas detectado com 6 chamadas.}
\evid{evidencias/ddg/fingerprinting-canvas/fingerprinting-canvas-crop.png}{1}{Canvas visual check e métodos interceptados.}

\subsection{Tracker Blocking (Request Blocking)}

\begin{table}[H]
\centering\small
\begin{tabular}{|p{2.4cm}|p{3.2cm}|p{3.6cm}|p{4.8cm}|}
\hline
\textbf{Teste} & \textbf{Esperado (página)} & \textbf{Plugin} & \textbf{Explicação} \\ \hline
Request Blocking & A página tenta 23 tipos de requisição a \texttt{bad.\allowbreak{}third-party.\allowbreak{}site} e orienta adicionar esse domínio à lista de bloqueio da ferramenta testada; sem bloqueio, todas carregam & 22 de 23 carregaram (apenas WebSocket falhou); plugin registrou \texttt{third-party.\allowbreak{}site}, 2 indicadores de hijacking (WebSocket e polling para \texttt{third-party.\allowbreak{}site}); score 85 & Por padrão o plugin apenas observa e só bloqueia domínios da lista personalizada, que estava vazia neste teste. Por isso o resultado coincide com um navegador sem bloqueio. A falha do WebSocket ocorreu na conexão, e o plugin registrou a tentativa em \texttt{onBeforeRequest}. O indicador de polling é compatível com as reconexões automáticas do teste de \texttt{server-sent-events}, que repete requisições ao mesmo terceiro. É um exemplo de comportamento que o plugin sinaliza sem conseguir distinguir intenção maliciosa. \\ \hline
\end{tabular}
\caption{Resultado de Request Blocking.}
\end{table}

A capacidade de bloqueio é demonstrada separadamente na Seção~\ref{sec:blocklist}. O resultado completo exportado pela página está em \path{evidencias/ddg/request-blocking/request-blocking-results.json}.

\evid{evidencias/ddg/request-blocking/request-blocking-crop.png}{1}{Request Blocking: resultados da página e popup.}
\evid{evidencias/ddg/request-blocking/request-blocking-1-crop.png}{1}{Request Blocking: indicadores de WebSocket e polling.}

\subsection{Storage Partitioning}

\begin{table}[H]
\centering\small
\begin{tabular}{|p{2.4cm}|p{3.2cm}|p{3.6cm}|p{4.8cm}|}
\hline
\textbf{Teste} & \textbf{Esperado (página)} & \textbf{Plugin} & \textbf{Explicação} \\ \hline
Storage Partitioning & Dados de 21 mecanismos isolados entre origens; resultado: 19 \emph{pass} e 2 \emph{unsupported} (WebSQL e Prefetch Cache) & Em \texttt{first-party.\allowbreak{}site}: 1 terceiro (\texttt{privacy-test-pages.\allowbreak{}site}), 1 cookie de primeira parte persistente; score 99 & O particionamento é aplicado pelo próprio Firefox (Total Cookie Protection e \emph{dynamic state partitioning}), e não pelo plugin. O plugin observa o armazenamento apenas do documento principal e não verifica se os dados de um iframe de terceiro estão particionados, então não há divergência a corrigir, mas também não há detecção. Em outra página visitada durante a coleta, o console do Firefox registrou uma mensagem de acesso particionado concedido a um embed do YouTube, que é esse mesmo mecanismo em um site real. \\ \hline
\end{tabular}
\caption{Resultado de Storage Partitioning.}
\end{table}

\evid{evidencias/ddg/storage-partitioning/storage-partitioning-1-crop.png}{1}{Storage Partitioning: resultados e popup.}
\evid{evidencias/ddg/storage-partitioning/storage-partitioning-2-crop.png}{1}{Storage Partitioning: continuação do popup.}

\subsection{Bounce Tracking}

\begin{table}[H]
\centering\small
\begin{tabular}{|p{2.4cm}|p{3.2cm}|p{3.6cm}|p{4.8cm}|}
\hline
\textbf{Teste} & \textbf{Esperado (página)} & \textbf{Plugin} & \textbf{Explicação} \\ \hline
Bounce Tracking (\emph{Go to good.third-party.site}) & \texttt{bad.\allowbreak{}third-party.\allowbreak{}site} lê seu UID e o repassa ao destino via URL; a página exibiu o ID 54 no \texttt{localStorage} e no cookie & 1 indicador de bounce tracking: \texttt{good.\allowbreak{}third-party.\allowbreak{}site} $\rightarrow$ \texttt{bad.\allowbreak{}third-party.\allowbreak{}site} $\rightarrow$ \texttt{good.\allowbreak{}third-party.\allowbreak{}site}, com 0,4 s no site intermediário & Durante o salto, o \texttt{bad.\allowbreak{}third-party.\allowbreak{}site} é o site principal da aba e acessa seu próprio armazenamento como primeira parte, contornando o bloqueio de cookies de terceiros. O plugin detecta o padrão pela permanência de 0,4 s, inferior ao limite de 2 s. Como origem, intermediário e destino compartilham o eTLD+1 \texttt{third-party.\allowbreak{}site}, a detecção também exigiu os parâmetros \texttt{bounceUID*} na URL de destino, o que evita confundir o bounce com um redirecionamento comum entre subdomínios. O ID permaneceu 54 entre execuções, o que mostra que o rastreador reconheceu o usuário; uma proteção completa apagaria o armazenamento do site intermediário, o que o plugin não faz. \\ \hline
\end{tabular}
\caption{Resultado de Bounce Tracking.}
\end{table}

\evid{evidencias/ddg/bounce-tracking/bounce-tracking-crop.png}{1}{Bounce Tracking: página de destino com o UID na URL e a cadeia detectada no popup.}

\subsection{Query Parameters}

\begin{table}[H]
\centering\small
\begin{tabular}{|p{2.4cm}|p{3.2cm}|p{3.6cm}|p{4.8cm}|}
\hline
\textbf{Teste} & \textbf{Esperado (página)} & \textbf{Plugin} & \textbf{Explicação} \\ \hline
Link com \texttt{fbclid}, \texttt{fb\_source} e 1 parâmetro comum & Resultado esperado \texttt{"u=14"} & URL final \texttt{query.\allowbreak{}html?u=14}; nenhum parâmetro de rastreamento sinalizado & O Firefox remove parâmetros de rastreamento conhecidos como \texttt{fbclid} antes da navegação (\emph{query stripping}, ativo em janela privada). A página recebeu exatamente o resultado esperado, e o plugin, que analisa a URL efetivamente carregada, não tinha o que sinalizar. Parâmetros que o Firefox não remove, como \texttt{utm\_*}, são sinalizados pelo plugin. \\ \hline
\end{tabular}
\caption{Resultado de Query Parameters.}
\end{table}

\evid{evidencias/ddg/query-parameters/query-parameters-crop.png}{1}{Query Parameters: URL já sem \texttt{fbclid}.}

\subsection{js-leaks e indicadores de hook}

A página \texttt{security/\allowbreak{}js-leaks.\allowbreak{}html} compara as propriedades do objeto \texttt{window} com o perfil de referência do Firefox 92 e lista o que foi adicionado, removido ou alterado. O teste foi executado com o plugin ligado e desligado.

\begin{table}[H]
\centering\small
\begin{tabular}{|p{2.4cm}|p{3.2cm}|p{3.6cm}|p{4.8cm}|}
\hline
\textbf{Teste} & \textbf{Esperado (página)} & \textbf{Plugin} & \textbf{Explicação} \\ \hline
js-leaks & Lista de diferenças no escopo global em relação ao perfil de referência & Com o plugin: \texttt{window.\allowbreak{}open} alterado, com código de \texttt{monitoredOpen}. Sem o plugin: \texttt{window.\allowbreak{}open} não aparece. Popup: 0 indicadores de hijacking & A página revela o rastro do próprio plugin, que envolve \texttt{window.\allowbreak{}open} para detectar pop-ups. O \texttt{hook-detector.\allowbreak{}js} ignora a própria substituição, por isso o popup não acusa hook. Os hooks de canvas não aparecem porque a página inspeciona as propriedades de \texttt{window}, e não os protótipos. As demais diferenças (APIs novas do Firefox 156, idiomas e \texttt{location.\allowbreak{}toString}) aparecem nas duas execuções e decorrem da versão e do perfil. \texttt{clipboard.\allowbreak{}writeText} alterado nas duas execuções corresponde ao scriptlet \texttt{prevent-clipboard-write} do uBlock Origin, que ainda estava ativo neste teste, e mostra que bloqueadores também deixam rastro observável. \\ \hline
\end{tabular}
\caption{Resultado de js-leaks.}
\end{table}

Esse resultado mostra um compromisso de projeto. O mesmo mecanismo que o BeEF usa para instrumentar uma página, a substituição de funções globais, é usado pelo plugin para detectar comportamentos e pelo uBlock Origin para neutralizar anúncios, e em todos os casos a alteração pode ser observada pela página, que pode usá-la para identificar a extensão ou para evitá-la.

\evid{evidencias/ddg/js-leaks/js-leaks-plugin-ligado-crop.png}{1}{js-leaks com o plugin ligado: \texttt{window.\allowbreak{}open} alterado.}
\evid{evidencias/ddg/js-leaks/js-leaks-plugin-desligado-crop.png}{1}{js-leaks com o plugin desligado: \texttt{window.\allowbreak{}open} não aparece.}

Como a página js-leaks não executa nenhum hook, a detecção foi validada também com um teste controlado em \texttt{example.\allowbreak{}com}, uma página sem scripts de terceiros: o comando \texttt{window.fetch = () => \{\}} executado no console gerou exatamente um indicador, \emph{Objeto global sobrescrito: window.fetch}, sem cookies, terceiros ou outros indicadores.

\evid{evidencias/ddg/hijacking/hijacking-fetch-crop.png}{1}{Teste controlado de hook em \texttt{example.\allowbreak{}com}.}

\subsection{Lista de bloqueio personalizada}\label{sec:blocklist}

Com \texttt{doubleclick.\allowbreak{}net} adicionado à lista e o uBlock Origin desativado, a página \emph{Major tracker via script} foi recarregada. O popup registrou 1 requisição bloqueada para \texttt{doubleclick.\allowbreak{}net} e exibiu o domínio na lista, com a opção de removê-lo. O domínio foi removido da lista antes das demais coletas.

\evid{evidencias/ddg/blocked/blocked-crop.png}{1}{Lista de bloqueio: requisição para \texttt{doubleclick.\allowbreak{}net} cancelada pelo plugin.}

\section{Análise de sites reais}

Conforme orientação do professor, os sites foram escolhidos livremente. Foram escolhidos o G1 (portal de notícias financiado por publicidade), o site institucional da UBC e o YouTube (plataforma de vídeo), por representarem perfis distintos de uso de terceiros e armazenamento.

Para cada site há, em \path{evidencias/sites-reais/<site>/}: capturas do popup (\path{privacy-inspector*.png}), o HAR sanitizado (\path{<site>.sanitized.har}), o resultado do Blacklight (\texttt{blacklight*.\allowbreak{}png}) e o popup e o logger do uBlock Origin (\texttt{ublock.\allowbreak{}png}, \texttt{ublock-logger.\allowbreak{}png} e o registro completo em \texttt{ublock-logger.\allowbreak{}md}).

\begin{table}[H]
\centering\small
\begin{tabular}{|l|c|c|c|c|c|c|}
\hline
\textbf{Site} & \textbf{Req. HAR} & \textbf{1ª parte} & \textbf{3ª parte} & \textbf{Status 0} & \textbf{Terceiros (plugin)} & \textbf{Score} \\ \hline
G1 & 181 & 44 & 137 & 19 & 21 & 62 (B) \\ \hline
UBC & 38 & 19 & 19 & 9 & 8 & 96 (A) \\ \hline
YouTube & 51 & 35 & 16 & 4 & 7 & 90 (A) \\ \hline
\end{tabular}
\caption{Resumo das coletas de 28/09. ``Status 0'' são requisições sem resposta no HAR, em sua maioria bloqueadas pelo ETP; as demais são beacons interrompidos, como os de \texttt{log\_event} do YouTube.}
\end{table}

\begin{table}[H]
\centering\small
\begin{tabular}{|p{3.4cm}|p{3.6cm}|p{3.6cm}|p{3.6cm}|}
\hline
\textbf{Ferramenta} & \textbf{G1} & \textbf{UBC} & \textbf{YouTube} \\ \hline
Blacklight: ad trackers & 67 & 12 & 0 (28/09); 3 (26/09) \\ \hline
Blacklight: cookies de terceiros & 100 & 15 & 0 \\ \hline
Blacklight: outros achados & Pixel do X; remarketing do Google Analytics & Pixel do X; remarketing do Google Analytics & Nenhum \\ \hline
uBlock Origin: bloqueios & 50 (21\%); 6 de 22 domínios conectados & 1 (2\%); 5 de 6 domínios & 5 (8\%); 4 de 6 domínios \\ \hline
Plugin: cookies & 13 de 1ª parte persistentes (subindo para 31 com a página aberta); 0 de 3ª & 0 & 2 de 1ª e 2 de 3ª, todos persistentes \\ \hline
Plugin: HTML5 & 11 local, 6 session, 0 IDB & 0 & 2 local, 0 session, 5 IDB \\ \hline
Plugin: canvas / sync & não / 0 & não / 0 & não / 0 \\ \hline
Plugin: hijacking & 4 (\texttt{fetch}, \texttt{XHR.\allowbreak{}open}, \texttt{XHR.\allowbreak{}send}, \texttt{addEventListener}) & 0 & 0 \\ \hline
\end{tabular}
\caption{Comparação entre o plugin, o Blacklight e o uBlock Origin.}
\end{table}

\subsection{G1}

O plugin registrou 21 domínios de terceiros e o HAR, 20. A diferença é \texttt{google-analytics.\allowbreak{}com}, que o plugin observou depois da exportação do HAR, já que a página continua fazendo requisições enquanto está aberta. O HAR registra 33 cabeçalhos \texttt{Set-Cookie}, todos de primeira parte e correspondentes a três cookies (\texttt{glb\_uid}, \texttt{glb\_uid\_jwt} e \texttt{hsid}), cada um redefinido 11 vezes. Isso explica o crescimento da contagem do plugin de 13 para 27 e 31 definições entre as capturas.

\small
\begin{longtable}{|p{3.3cm}|p{1.9cm}|p{1.6cm}|p{1.6cm}|p{5.0cm}|}
\caption{Reconciliação de domínios no G1.}\\
\hline
\textbf{Domínio} & \textbf{Plugin / HAR} & \textbf{Blacklight} & \textbf{uBO} & \textbf{Explicação} \\ \hline
\endfirsthead
\hline
\textbf{Domínio} & \textbf{Plugin / HAR} & \textbf{Blacklight} & \textbf{uBO} & \textbf{Explicação} \\ \hline
\endhead
\texttt{doubleclick.\allowbreak{}net} & sim (6 req., todas bloqueadas pelo ETP) & sim & bloqueado & Concordância. O plugin conta a tentativa de conexão, mesmo quando o ETP a cancela. \\ \hline
\texttt{googletagmanager.\allowbreak{}com}, \texttt{google.\allowbreak{}com}, \texttt{google-analytics.\allowbreak{}com} & sim & sim & bloqueado & Concordância. O uBO bloqueou \texttt{/\allowbreak{}ccm/\allowbreak{}collect}, \texttt{/\allowbreak{}rmkt/\allowbreak{}collect} e \texttt{/\allowbreak{}g/\allowbreak{}collect}; o HAR mostra \texttt{analytics.\allowbreak{}google.\allowbreak{}com} bloqueado pelo ETP. \\ \hline
\texttt{googlesyndication.\allowbreak{}com}, \texttt{rubiconproject.\allowbreak{}com}, \texttt{scorecardresearch.\allowbreak{}com} & sim (bloqueados pelo ETP) & sim & bloqueado & Concordância. O uBO substituiu \texttt{adsbygoogle.\allowbreak{}js} por um surrogate. \\ \hline
\texttt{ads-twitter.\allowbreak{}com} & sim (bloqueado) & sim (pixel do X) & bloqueado & Concordância. \\ \hline
\texttt{facebook.\allowbreak{}net} & sim (bloqueado) & não & bloqueado & O HAR mostra a tentativa de carregar \texttt{fbevents.\allowbreak{}js} na coleta feita no Brasil. O Blacklight, a partir da Califórnia, não encontrou o pixel, o que é compatível com tags configuradas por região. \\ \hline
\texttt{clarity.\allowbreak{}ms} & sim (bloqueado) & não (sem gravação de sessão) & bloqueado & O Microsoft Clarity é um serviço de gravação de sessão. Ele foi requisitado na coleta brasileira e não apareceu no scan do Blacklight, pela mesma razão. \\ \hline
\texttt{permutive.\allowbreak{}app}, \texttt{adnami.\allowbreak{}io}, \texttt{dv.\allowbreak{}tech}, \texttt{mrf.\allowbreak{}io}, \texttt{debugbear.\allowbreak{}com} & sim & não & bloqueado & Não constam entre as empresas listadas pelo Blacklight, que classifica rastreadores pela base Tracker Radar do DuckDuckGo de março de 2024. O uBO bloqueia pelas suas listas de filtros. \\ \hline
\texttt{g.\allowbreak{}globo}, \texttt{glbimg.\allowbreak{}com}, \texttt{globoi.\allowbreak{}com} & sim & não & parcial & São domínios do próprio Grupo Globo com eTLD+1 diferente de \texttt{globo.\allowbreak{}com}, por isso o plugin os conta como terceiros. O uBO bloqueou recursos específicos: beacons de \texttt{sdk-metrics.\allowbreak{}g.\allowbreak{}globo}, \texttt{sentry.\allowbreak{}globoi.\allowbreak{}com} e o Chartbeat em \texttt{s3.\allowbreak{}glbimg.\allowbreak{}com}. \\ \hline
\texttt{horizon-track.\allowbreak{}globo.\allowbreak{}com}, \texttt{trackid.\allowbreak{}globoid.\allowbreak{}globo.\allowbreak{}com} & não (1ª parte) & não & bloqueado & Rastreadores de primeira parte. Só o uBO os identifica, por filtrar URLs e não domínios. \\ \hline
\texttt{w3schools.\allowbreak{}com} & sim & não & não & Vídeo de exemplo (\texttt{/\allowbreak{}tags/\allowbreak{}movie.\allowbreak{}mp4}) embutido na página, sem relação com rastreamento; nenhuma ferramenta o bloqueia. \\ \hline
Criteo, PubMatic, OpenX, Amazon, BidSwitch, MediaMath, Lotame, TowerData, Adobe e outros (20 domínios) & não & sim & não & Esses domínios participam do leilão de anúncios (header bidding) e da sincronização de IDs. O leilão é disparado pelo Prebid (\path{ads.rubiconproject.com/prebid/11366_g1.js}), que foi bloqueado pelo ETP na coleta do plugin e pelo uBO. Sem o leilão, esses parceiros nunca são contatados, os 100 cookies de terceiros vistos pelo Blacklight não são criados e não ocorre o cookie sync que o plugin detectaria. \\ \hline
\end{longtable}
\normalsize

Os quatro indicadores de hook no G1 apareceram com o uBlock Origin desativado, portanto vêm de scripts da página. A combinação de \texttt{fetch}, \texttt{XMLHttpRequest.\allowbreak{}open/\allowbreak{}send} e \texttt{addEventListener} é a instrumentação típica de ferramentas de monitoramento de erros e desempenho, e o HAR mostra requisições para o Sentry (\texttt{sentry.\allowbreak{}globoi.\allowbreak{}com}) e para o DebugBear. É um caso de uso legítimo do mesmo mecanismo que um hook malicioso usaria, e o plugin não tem como diferenciar os dois.

\popups{evidencias/sites-reais/g1/privacy-inspector.png}{evidencias/sites-reais/g1/privacy-inspector-1.png}{evidencias/sites-reais/g1/privacy-inspector-2.png}{Popup do plugin no G1: score, domínios e armazenamento.}
\pair{evidencias/sites-reais/g1/privacy-inspector-3.png}{evidencias/sites-reais/g1/privacy-inspector-4.png}{Popup do plugin no G1: hijacking e cookies.}
\evid{evidencias/sites-reais/g1/blacklight.png}{0.6}{Blacklight no G1 (detalhes em \texttt{blacklight-detail.\allowbreak{}png}).}
\evid{evidencias/sites-reais/g1/ublock.png}{1}{uBlock Origin no G1: 50 bloqueios.}
\evid{evidencias/sites-reais/g1/ublock-logger.png}{0.8}{Logger do uBlock Origin no G1 (amostra; registro completo em \texttt{ublock-logger.\allowbreak{}md}).}

\subsection{UBC}

Plugin e HAR coincidem nos 8 domínios de terceiros. Nenhum cookie foi definido por cabeçalho \texttt{Set-Cookie} e nenhum armazenamento HTML5 foi usado. O Google Analytics, que foi executado, cria seus cookies por JavaScript (\texttt{document.\allowbreak{}cookie}), que o plugin não contabiliza.

\small
\begin{longtable}{|p{3.3cm}|p{1.9cm}|p{1.6cm}|p{1.6cm}|p{5.0cm}|}
\caption{Reconciliação de domínios na UBC.}\\
\hline
\textbf{Domínio} & \textbf{Plugin / HAR} & \textbf{Blacklight} & \textbf{uBO} & \textbf{Explicação} \\ \hline
\endfirsthead
\hline
\textbf{Domínio} & \textbf{Plugin / HAR} & \textbf{Blacklight} & \textbf{uBO} & \textbf{Explicação} \\ \hline
\endhead
\texttt{googletagmanager.\allowbreak{}com} & sim & sim & bloqueado (surrogate) & O uBO fez um único bloqueio: substituiu o \texttt{gtm.\allowbreak{}js} por um surrogate. Como o GTM carrega as demais tags, esse bloqueio impede todas as outras requisições de rastreamento, o que explica o uBO registrar só 1 bloqueio enquanto o Blacklight encontrou 12 rastreadores. \\ \hline
\texttt{doubleclick.\allowbreak{}net} (\texttt{ad.}, \texttt{fls.}) & sim (5 req., todas bloqueadas pelo ETP) & sim & não chegou a ser requisitado & Tags do Floodlight disparadas pelo GTM. O ETP as cancelou, então nenhum cookie foi criado. \\ \hline
\texttt{google.\allowbreak{}com}, \texttt{google.\allowbreak{}com.\allowbreak{}br} & sim & sim (\texttt{google.\allowbreak{}com}) & não chegou a ser requisitado & Google Analytics e conversões do Google Ads (\texttt{/\allowbreak{}g/\allowbreak{}collect}, \texttt{/\allowbreak{}ccm/\allowbreak{}collect}, \texttt{/\allowbreak{}gmp/\allowbreak{}conversion}), que não foram bloqueados. A requisição \path{google.com.br/ads/ga-audiences} corresponde ao recurso de \emph{remarketing audiences} do Google Analytics, o mesmo que o Blacklight apontou, e mostra concordância entre as ferramentas. \\ \hline
\texttt{sc-static.\allowbreak{}net}, \texttt{mookie1.\allowbreak{}com} & sim (bloqueados pelo ETP) & não listados individualmente & não chegou a ser requisitado & Pixel do Snapchat e Xaxis/GroupM. O Blacklight cita ``quatro outras empresas'' sem detalhá-las. \\ \hline
\texttt{ajax.\allowbreak{}googleapis.\allowbreak{}com}, \texttt{typography.\allowbreak{}com} & sim & não & não & Biblioteca jQuery e fontes; não são rastreadores. O plugin as conta por serem de outro domínio. \\ \hline
\texttt{linkedin.\allowbreak{}com}, \texttt{licdn.\allowbreak{}com}, X e Microsoft & não & sim & não & Tags disparadas pelo GTM que não foram requisitadas na coleta brasileira. É compatível com tags condicionadas por região ou disparadas em cadeia após tags que o ETP cancelou. \\ \hline
\end{longtable}
\normalsize

\popups{evidencias/sites-reais/ubc/privacy-inspector.png}{evidencias/sites-reais/ubc/privacy-inspector-1.png}{evidencias/sites-reais/ubc/privacy-inspector-2.png}{Popup do plugin na UBC.}
\evid{evidencias/sites-reais/ubc/blacklight.png}{0.6}{Blacklight na UBC (detalhes em \texttt{blacklight-detalhes.\allowbreak{}png}).}
\evid{evidencias/sites-reais/ubc/ublock.png}{1}{uBlock Origin na UBC: 1 bloqueio.}
\evid{evidencias/sites-reais/ubc/ublock-logger.png}{0.8}{Logger do uBlock Origin na UBC: \texttt{gtm.\allowbreak{}js} substituído por surrogate.}

\subsection{YouTube}

Plugin e HAR coincidem nos 7 domínios de terceiros. O plugin contabilizou 4 cookies persistentes (2 de primeira e 2 de terceira parte), mas o HAR exportado não contém cabeçalhos \texttt{Set-Cookie} para o YouTube. Uma hipótese é que parte das respostas tenha sido atendida pelo service worker do YouTube, cujos cabeçalhos não são incluídos na exportação da mesma forma. O plugin lê os cabeçalhos diretamente da API \texttt{webRequest}.

\small
\begin{longtable}{|p{3.3cm}|p{1.9cm}|p{1.6cm}|p{1.6cm}|p{5.0cm}|}
\caption{Reconciliação de domínios no YouTube.}\\
\hline
\textbf{Domínio} & \textbf{Plugin / HAR} & \textbf{Blacklight} & \textbf{uBO} & \textbf{Explicação} \\ \hline
\endfirsthead
\hline
\textbf{Domínio} & \textbf{Plugin / HAR} & \textbf{Blacklight} & \textbf{uBO} & \textbf{Explicação} \\ \hline
\endhead
\texttt{doubleclick.\allowbreak{}net} & sim (2 req., bloqueadas pelo ETP) & só em 26/09 & bloqueado & \path{static.doubleclick.net/instream/ad_status.js} foi trocado por surrogate pelo uBO, e \path{googleads.g.doubleclick.net/pagead/id}, que busca um ID de publicidade, recebeu resposta vazia (\texttt{noop.\allowbreak{}txt}). \\ \hline
\texttt{google.\allowbreak{}com}, \texttt{google.\allowbreak{}com.\allowbreak{}br} & sim & só em 26/09 & bloqueado & Requisições \texttt{/\allowbreak{}pagead/\allowbreak{}lvz} de conversão de anúncio. O \texttt{.\allowbreak{}com.\allowbreak{}br} é específico da coleta brasileira. \\ \hline
\texttt{ytimg.\allowbreak{}com} & sim & sim & não & Imagens do YouTube, único domínio listado pelo Blacklight em 28/09. \\ \hline
\texttt{googlevideo.\allowbreak{}com}, \texttt{gstatic.\allowbreak{}com}, \texttt{fonts.\allowbreak{}googleapis.\allowbreak{}com} & sim & não & não & CDN de vídeo, recursos estáticos e fontes. São da mesma organização (Alphabet), mas de outro domínio. \\ \hline
\path{youtube.com/youtubei/v1/log_event}, \texttt{/\allowbreak{}generate\_204} & não (1ª parte) & não & bloqueado & Telemetria de primeira parte, bloqueada apenas pelo uBO. \\ \hline
\texttt{google-analytics.\allowbreak{}com}, \texttt{googletagmanager.\allowbreak{}com} & não & só em 26/09 & não & Presentes apenas no scan de 26/09. O Blacklight passou de 3 para 0 rastreadores entre 26/09 e 28/09, o que mostra que a mesma página varia entre carregamentos. \\ \hline
\end{longtable}
\normalsize

\popups{evidencias/sites-reais/youtube/privacy-inspector.png}{evidencias/sites-reais/youtube/privacy-inspector-1.png}{evidencias/sites-reais/youtube/privacy-inspector-2.png}{Popup do plugin no YouTube.}
\evid{evidencias/sites-reais/youtube/blacklight.png}{0.6}{Blacklight no YouTube em 28/09 (o scan de 26/09 está em \texttt{blacklight-2026-09-26.\allowbreak{}png}).}
\evid{evidencias/sites-reais/youtube/ublock.png}{1}{uBlock Origin no YouTube: 5 bloqueios.}
\evid{evidencias/sites-reais/youtube/ublock-logger.png}{0.8}{Logger do uBlock Origin no YouTube: bloqueios de rede e scriptlets que removem anúncios do player.}

\section{Pontuação de privacidade}

\subsection{Metodologia}

A pontuação vai de 0 a 100, e quanto maior, melhor. Cada critério tem um peso, e os pesos somam 100, além de um limite de saturação. A penalidade de cada critério é proporcional ao valor observado até atingir o limite:
\[
\text{penalidade}_i = \text{peso}_i \times \min\left(\frac{\text{valor}_i}{\text{limite}_i}, 1\right),
\qquad
\text{score} = 100 - \sum_i \text{penalidade}_i .
\]

O critério usado para definir os pesos foi o quanto cada mecanismo permite que terceiros identifiquem o usuário entre sites e compartilhem ou comercializem esses dados, ou ajam sem o seu consentimento.

\begin{table}[H]
\centering\small
\begin{tabular}{|p{3.4cm}|c|c|p{7.6cm}|}
\hline
\textbf{Critério} & \textbf{Peso} & \textbf{Limite} & \textbf{Justificativa} \\ \hline
Bounce tracking / cookie sync & 25 & 3 & É a troca direta do identificador do usuário entre empresas, ou seja, o compartilhamento de dados acontecendo. \\ \hline
Cookies de terceira parte & 20 & 10 & Identificador clássico que acompanha o usuário entre sites. \\ \hline
Indicadores de hijacking/hook & 20 & 3 & Terceiros executando ações sem consentimento (abrir abas, redirecionar, interceptar requisições), com possibilidade de exfiltrar dados. Inclui pop-unders, como abas de apostas abertas automaticamente. \\ \hline
Canvas fingerprinting & 15 & 1 & Identifica o usuário entre sites e não pode ser apagado por ele. O peso é menor que o de cookies de terceiros porque o detector conta qualquer uso de canvas, com risco de falso positivo. \\ \hline
Domínios de terceiros & 10 & 20 & Cada conexão vaza o IP e a página visitada, mas sozinha não identifica o usuário, e muitos domínios são CDNs ou da mesma organização. \\ \hline
Cookies persistentes & 5 & 20 & Em geral são da própria página e indicam duração, não compartilhamento. \\ \hline
Armazenamento HTML5 & 5 & 20 & Medido apenas no documento principal, ou seja, dados da própria página. \\ \hline
\end{tabular}
\caption{Critérios, pesos e justificativas.}
\end{table}

Classificação: A (80 a 100, boa privacidade), B (60 a 79, moderada), C (40 a 59, baixa) e D (abaixo de 40, muito baixa). O cálculo está em \texttt{score.\allowbreak{}js} e é exibido no popup com a penalidade de cada critério.

\subsection{Aplicação aos três sites}

\begin{table}[H]
\centering\small
\begin{tabular}{|l|c|c|c|}
\hline
\textbf{Critério (peso)} & \textbf{G1} & \textbf{UBC} & \textbf{YouTube} \\ \hline
Domínios de terceiros (10) & 21/20 $\rightarrow$ $-10$ & 8/20 $\rightarrow$ $-4$ & 7/20 $\rightarrow$ $-3{,}5$ \\ \hline
Cookies de terceira parte (20) & 0/10 $\rightarrow$ 0 & 0/10 $\rightarrow$ 0 & 2/10 $\rightarrow$ $-4$ \\ \hline
Cookies persistentes (5) & 13/20 $\rightarrow$ $-3{,}3$ & 0/20 $\rightarrow$ 0 & 4/20 $\rightarrow$ $-1$ \\ \hline
Armazenamento HTML5 (5) & 17/20 $\rightarrow$ $-4{,}3$ & 0/20 $\rightarrow$ 0 & 7/20 $\rightarrow$ $-1{,}8$ \\ \hline
Canvas fingerprinting (15) & 0 & 0 & 0 \\ \hline
Bounce / cookie sync (25) & 0 & 0 & 0 \\ \hline
Hijacking/hook (20) & 4/3 $\rightarrow$ $-20$ & 0 & 0 \\ \hline
\textbf{Score} & \textbf{62 (B)} & \textbf{96 (A)} & \textbf{90 (A)} \\ \hline
Blacklight (ad trackers / cookies 3ª) & 67 / 100 & 12 / 15 & 0 / 0 \\ \hline
\end{tabular}
\caption{Pontuação dos sites reais e resultado do Blacklight.}
\end{table}

\subsection{Comparação crítica com o Blacklight}

\textbf{Onde concordam.} As duas ferramentas apontam o G1 como o site mais invasivo dos três, e o YouTube como um site com pouco rastreamento de terceiros na página inicial sem login. Os rastreadores que o plugin viu no G1 e que são reconhecidos como tal (DoubleClick, Google, Rubicon, Comscore, pixel do X) também foram listados pelo Blacklight.

\textbf{Onde divergem e por quê.}
\begin{enumerate}
  \item \textbf{O que é medido.} O Blacklight mede o que o site \emph{tenta fazer} com um visitante sem nenhuma proteção. O plugin mede o que \emph{efetivamente aconteceu} no navegador da coleta, que tinha o ETP ativo. No G1, o ETP bloqueou o Prebid, e com isso o leilão de anúncios não ocorreu. Por isso o plugin registrou 0 cookies de terceiros contra 100 do Blacklight, e nenhum cookie sync, embora o Blacklight tenha visto parceiros típicos de sincronização como BidSwitch, TowerData, Lotame e Adobe.
  \item \textbf{A inversão entre UBC e YouTube.} No plugin, a UBC tem a melhor nota (96): as tags de DoubleClick/Floodlight, Snapchat e Xaxis disparadas pelo GTM foram bloqueadas pelo ETP, e as que passaram (Google Analytics e conversões do Google) não definiram cookies por cabeçalho, já que o Google Analytics cria os seus por JavaScript. No Blacklight, a UBC tem 12 rastreadores e 15 cookies de terceiros (X, LinkedIn, Microsoft), enquanto o YouTube tem zero. O YouTube perdeu pontos no plugin por 2 cookies de \texttt{google.\allowbreak{}com} e por domínios como \texttt{googlevideo.\allowbreak{}com}, que são da mesma organização. Para o plugin, domínio diferente é terceiro; para o Blacklight, só conta o que está na lista de rastreadores.
  \item \textbf{Hooks.} O Blacklight não avalia hooks em objetos globais. O plugin aplicou a penalidade máxima do critério ao G1 por quatro funções sobrescritas, que muito provavelmente vêm de ferramentas de monitoramento como o Sentry. Isso reduz a nota do G1 por um comportamento que não é necessariamente malicioso, e é a principal fonte de falso positivo da metodologia.
  \item \textbf{Localização e tempo.} O Blacklight escaneia a partir da Califórnia, e a coleta do plugin foi feita no Brasil: Facebook e Microsoft Clarity só apareceram na coleta brasileira, LinkedIn e X só no scan americano, e \texttt{google.\allowbreak{}com.\allowbreak{}br} só no Brasil. Além disso, a mesma página varia entre carregamentos: o YouTube teve 3 rastreadores no Blacklight em 26/09 e 0 em 28/09, e a contagem de cookies do G1 subiu de 13 para 31 enquanto a página permanecia aberta.
  \item \textbf{Saturação.} Com limite de 20 domínios, o G1 atinge a penalidade máxima desse critério, e qualquer domínio adicional deixa de pesar. O mesmo vale para os três indicadores de hijacking. É uma escolha deliberada para que um único critério não domine a nota, mas ela reduz a capacidade de distinguir sites muito invasivos entre si.
\end{enumerate}

Em síntese, o score do plugin é uma medida da exposição efetiva do usuário no navegador em que foi coletado, enquanto o Blacklight é uma medida do comportamento potencial do site. As duas são complementares, e as divergências ficam explicáveis pela configuração de proteção, pela localização e pela diferença entre ``domínio de terceiro'' e ``rastreador conhecido''.

\section{Limitações}

\begin{itemize}
  \item Domínio de terceiro não é sinônimo de rastreador. O plugin não usa listas de rastreadores e conta como terceiros CDNs e domínios da mesma organização.
  \item Cookies definidos por JavaScript não são contabilizados, e redefinições do mesmo cookie contam várias vezes.
  \item O armazenamento HTML5 é medido apenas no documento principal e em um único momento, cerca de 1 segundo após o carregamento.
  \item O detector de canvas conta qualquer leitura de canvas, sem distinguir uso legítimo de fingerprinting, e não monitora WebGL, fontes, \texttt{navigator} ou \texttt{screen}.
  \item Os indicadores de hook não distinguem instrumentação legítima (monitoramento, bloqueadores) de código malicioso, e as substituições feitas pelo próprio plugin são observáveis pela página, como mostrou o teste js-leaks.
  \item A abertura de abas é sinalizada mesmo quando resulta de um clique legítimo, e fluxos legítimos de login (SSO/OAuth) podem ser sinalizados como bounce tracking.
  \item O plugin detecta, mas não impede, bounce tracking e cookie sync, e não apaga o armazenamento de sites intermediários.
  \item Os resultados dependem da configuração do navegador: com o ETP ativo, parte do rastreamento é bloqueada antes de ser observada.
\end{itemize}

\section{Conclusão}

O Privacy Inspector detecta conexões com terceiros, cookies por tipo, armazenamento HTML5, canvas fingerprinting, bounce tracking, indicadores de hijacking e hook e calcula uma pontuação de privacidade com metodologia explícita, além de permitir o bloqueio de domínios escolhidos pelo usuário. Nas DuckDuckGo Privacy Test Pages, o plugin detectou todos os rastreadores de Tracker Reporting, inclusive os carregados dinamicamente ou com atraso, as chamadas de canvas, o bounce tracking e os indicadores de WebSocket e polling. As divergências encontradas se explicam por proteções que o próprio Firefox já aplica (\emph{query stripping} e particionamento de armazenamento) ou por funcionalidades que o plugin não se propõe a oferecer, como surrogates.

Nos sites reais, a comparação com o HAR, o Blacklight e o uBlock Origin mostrou que cada ferramenta responde a uma pergunta diferente: o plugin observa o que aconteceu no navegador, o Blacklight revela o que o site tenta fazer sem proteção, e o uBlock Origin mostra o que uma lista de filtros impede, inclusive rastreadores de primeira parte. A principal lição da análise é que o bloqueio de um único script, como o Prebid no G1 ou o Google Tag Manager na UBC, elimina em cadeia dezenas de rastreadores. Por isso, contagens isoladas de domínios ou cookies precisam ser interpretadas junto com as condições da coleta.

\end{document}
